\section{Conclusion and Future Work}
\label{sec:conclusion}

This work introduces \benchname, a unified benchmark for systematically evaluating memory-augmented robotic manipulation across four cognitive dimensions: temporal, spatial, object, and procedural memory. We further develop a family of vision-language-action (VLA) models and conducted controlled comparisons of symbolic, perceptual, and recurrent memory representations with multiple integration mechanisms. Results show that no single design consistently dominates: symbolic memory excels at counting and short-horizon reasoning, while perceptual memory is crucial for motion-centric and time-sensitive behaviors.


Despite its breadth, \benchname\xspace focuses on tabletop manipulation with a fixed set of assets and evaluates a single pretrained backbone ($\pi_{0.5}$), leaving mobile manipulation and alternative architectures, such as memory-banks and other VLA backbones, for future study. 
Our results further suggest that memory representations are complementary, motivating unified frameworks that integrate multiple forms of memory. We hope \benchname\xspace serves as a foundation and lasting testbed for developing memory-augmented robotic policies.

%% Acknowledgments are kept here for the camera-ready version. Per CVPR policy,
%% they must be omitted from the blind-review submission to preserve anonymity.
\section*{Acknowledgments}
This work was supported in part by NSF IIS-1949634, NSF SES-2128623, NSF NAIRR250085, NSF CAREER \#2337870, and NSF NRI \#2220876. This research is also supported by the National Artificial Intelligence Research Resource (NAIRR) Pilot and the Delta advanced computing and data resource which is supported by the National Science Foundation (award NSF-OAC 2005572).
